% ============================================================================
\section{Implémentation distribuée et orchestration}
\label{sec:orchestration}
% ============================================================================

\subsection{Déploiement}

Les VM sont gérées par des scripts \code{gcloud}. Elles sont décrites à un seul
endroit, \code{data/machines.csv} (nom, type GCP, vCPU, RAM, prix, contention),
lu à la fois par ces scripts et par le scheduler : changer de parc de VM revient à
modifier ce fichier. Seuls le projet et la zone GCP sont dans
\code{scripts/config.sh}.

\begin{lstlisting}
./scripts/manage_vms.sh create   # crée les VM manquantes, corrige le type si besoin
./scripts/manage_vms.sh start
./scripts/deploy.sh              # compile en linux/amd64 et copie les binaires
./scripts/manage_vms.sh stop     # toujours à la fin d'une session
\end{lstlisting}

\code{deploy.sh} compile les quatre workloads \textbf{une seule fois}, sur la
machine de contrôle, pour Linux x86\_64 (\code{GOOS=linux GOARCH=amd64}, binaires
statiques), puis les copie sur chaque VM et vérifie qu'ils s'y exécutent. Aucune
installation n'est nécessaire sur les VM.

\subsection{L'orchestrateur}

L'orchestrateur (\code{orchestrator/main.go}) exécute une allocation telle quelle :
il ne déplace jamais une tâche vers une autre VM que celle choisie par le scheduler.

\begin{itemize}
  \item \textbf{Une file par VM}, dans l'ordre de \code{allocation.csv}. Toutes les
        VM travaillent en parallèle ; sur chaque VM, \textbf{N workers} (par défaut
        N = nombre de vCPU, section~\ref{sec:workers}) vident la file : quand un
        worker finit une tâche, il prend la suivante.
  \item \textbf{Une seule connexion SSH par VM}, ouverte au départ et réutilisée par
        toutes ses tâches (multiplexage OpenSSH \code{ControlMaster}), avec la clé et
        la vérification d'identité d'hôte de \code{gcloud}.
  \item \textbf{Horodatage sur la machine de contrôle} : début et fin de chaque
        tâche sont mesurés avec une seule horloge ; aucune synchronisation entre VM
        n'est nécessaire.
  \item \textbf{Séparation calcul / surcoût} : chaque workload renvoie son propre
        temps de calcul ; l'orchestrateur en déduit
        \code{overhead\_s = durée mesurée $-$ runtime\_s}, le coût du réseau et de SSH.
  \item \textbf{Robustesse} : une tâche en échec n'arrête pas les autres (code de
        retour et erreur enregistrés) ; au démarrage, l'orchestrateur attend que le
        serveur SSH des VM soit prêt (une VM \code{RUNNING} refuse encore les
        connexions pendant 20 à 40~s) ; il relève le modèle de processeur de chaque VM.
\end{itemize}

\paragraph{Sorties.} \code{results/<run>.csv} contient une ligne par tâche
(\code{task\_id, workload, size, vm, worker, start\_s, end\_s, duration\_s,
remote\_runtime\_s, overhead\_s, exit\_code, error}) ; \code{results/<run>\_vms.csv}
résume chaque VM (fenêtre active, temps occupé, surcoût moyen, durée facturée,
coût, processeur).

\subsection{Premières preuves d'exécution distribuée}

Premier lot réel : les 9 tâches du TD3 réparties par Round Robin sur les 4 VM.

\begin{table}[H]
\centering\small
\begin{tabular}{lrrrr}
\toprule
\textbf{VM} & \textbf{Tâches} & \textbf{Fenêtre active (s)} & \textbf{Temps de calcul (s)} & \textbf{Surcoût SSH / tâche} \\
\midrule
highcpu & 2 & 0,54 & 0,58 & 0,077 s \\
vm1     & 3 & 1,31 & 1,68 & 0,085 s \\
vm2     & 2 & 0,83 & 1,45 & 0,070 s \\
vm3     & 2 & 0,68 & 0,94 & 0,099 s \\
\bottomrule
\end{tabular}
\caption{Premier lot réel (\mesure{}) : makespan 1,316~s, 9 tâches sur 9 réussies.}
\end{table}

Depuis, l'orchestrateur a exécuté près de 40\,000 tâches sur GCP au cours des
expériences (sections~\ref{sec:workers} et~\ref{sec:reels}), sans perte de tâche.

\paragraph{Le surcoût SSH.} Grâce à la connexion réutilisée, il est de
\textbf{0,06 à 0,12~s par tâche}, contre environ 0,5 à 1~s avec une connexion par
tâche (ordre de grandeur de notre première version). Il augmente avec le nombre
de tâches simultanées (jusqu'à 0,19~s avec 4 workers sur une VM à 2~vCPU) : nous
retenons 0,1~s dans le modèle.
